% Valid Time Integration for Elastica in the Snake-HRL Project
% Standalone reference document — defines the two-level stepping hierarchy,
% stability constraints, and CPG–integrator synchronization

\documentclass[11pt]{article}

% ---------------------------------------------------------------------------
% Font and encoding
% ---------------------------------------------------------------------------
\usepackage{palatino}
\usepackage[T1]{fontenc}
\usepackage[utf8]{inputenc}

% ---------------------------------------------------------------------------
% Page layout
% ---------------------------------------------------------------------------
\usepackage[top=1in, bottom=1in, left=1.25in, right=1.25in]{geometry}
\usepackage{setspace}
\setstretch{1.15}
\setlength{\parindent}{1em}
\setlength{\parskip}{0.4em}

% ---------------------------------------------------------------------------
% Math
% ---------------------------------------------------------------------------
\usepackage{amsmath}
\usepackage{amssymb}
\usepackage{mathtools}
\usepackage{bm}

% ---------------------------------------------------------------------------
% Figures and tables
% ---------------------------------------------------------------------------
\usepackage{booktabs}
\usepackage{float}
\usepackage{array}
\usepackage{tabularx}

% ---------------------------------------------------------------------------
% Colors and hyperlinks
% ---------------------------------------------------------------------------
\usepackage{xcolor}
\usepackage[hidelinks]{hyperref}

% ---------------------------------------------------------------------------
% Lists
% ---------------------------------------------------------------------------
\usepackage{enumitem}
\setlist[itemize]{nosep, topsep=0.3em, itemsep=0.2em}

% ---------------------------------------------------------------------------
% Section formatting
% ---------------------------------------------------------------------------
\usepackage{titlesec}
\titlespacing*{\section}{0pt}{1.2em}{0.6em}
\titlespacing*{\subsection}{0pt}{0.9em}{0.3em}

% ---------------------------------------------------------------------------
% Algorithms
% ---------------------------------------------------------------------------
\usepackage{algorithm}
\usepackage{algpseudocode}

% ---------------------------------------------------------------------------
% Custom commands
% ---------------------------------------------------------------------------
\newcommand{\vect}[1]{\mathbf{#1}}
\newcommand{\mat}[1]{\mathbf{#1}}
\newcommand{\Real}{\mathbb{R}}
\newcommand{\ddt}[1]{\frac{\partial #1}{\partial t}}
\newcommand{\dds}[1]{\frac{\partial #1}{\partial s}}
\newcommand{\dd}[2]{\frac{d #1}{d #2}}
\newcommand{\dtrl}{\Delta t_{\mathrm{RL}}}
\newcommand{\dtsub}{\Delta t_{\mathrm{sub}}}
\newcommand{\Nsub}{N_{\mathrm{sub}}}
\newcommand{\Nrl}{N_{\mathrm{RL}}}
\newcommand{\krest}{\bm{\kappa}_{\mathrm{rest}}}

% ---------------------------------------------------------------------------
\begin{document}
% ---------------------------------------------------------------------------

\begin{center}
  {\Large\bfseries Valid Time Integration for PyElastica} \\[0.5em]
  {\large Two-Level Stepping, Stability Constraints, and CPG Synchronization} \\[1.5em]
  \textit{Snake-HRL Project --- Reference Document}
\end{center}

\vspace{1em}

% ===================================================================
\section{Overview}
\label{sec:overview}
% ===================================================================

PyElastica simulates Cosserat rod dynamics using explicit symplectic integrators.
In the Snake-HRL project, the simulation is embedded inside a reinforcement learning
(RL) loop, producing a \emph{two-level time-stepping hierarchy}: a fast inner loop
that integrates the rod mechanics, and a slower outer loop at which the RL agent
selects actions.  A \emph{valid} time integration is one in which:
\begin{enumerate}[nosep]
  \item The inner substep $\dtsub$ is small enough for the symplectic integrator
        to remain stable given the rod's material stiffness and geometry.
  \item The CPG (Central Pattern Generator) curvature field is updated at every
        substep, not just at RL boundaries, so the actuation is temporally smooth.
  \item Friction forces are applied before each integration substep so they
        participate in the force balance.
  \item Numerical damping is low enough to preserve locomotion dynamics but
        high enough to suppress spurious high-frequency modes.
\end{enumerate}

% ===================================================================
\section{Cosserat Rod Equations of Motion}
\label{sec:eom}
% ===================================================================

Let $\vect{r}(s,t) \in \Real^3$ denote the centerline position and
$\mat{Q}(s,t) \in SO(3)$ the material frame of a Cosserat rod parameterized
by arc length $s \in [0, L]$.  The governing PDEs are:

\begin{align}
  \rho A \, \ddot{\vect{r}}
    &= \dds{\vect{n}} + \vect{f}_{\mathrm{ext}},
    \label{eq:linear-momentum} \\[4pt]
  \rho \vect{I} \, \dot{\bm{\omega}}
    &= \dds{\vect{m}} + \dds{\vect{r}} \times \vect{n}
       + \bm{\tau}_{\mathrm{ext}},
    \label{eq:angular-momentum}
\end{align}

\noindent where $\vect{n}$ is the internal force resultant, $\vect{m}$ the
internal moment, $\rho$ the density, $A$ the cross-sectional area,
$\vect{I}$ the second moment of area tensor, $\bm{\omega}$ the angular
velocity, and $\vect{f}_{\mathrm{ext}}$, $\bm{\tau}_{\mathrm{ext}}$ are
external forces and torques (gravity, friction).

The constitutive relations for an isotropic rod are:
\begin{align}
  \vect{m} &= \mat{B}\,(\bm{\kappa} - \krest), \label{eq:bending} \\
  \vect{n} &= \mat{S}\,(\bm{\sigma} - \bm{\sigma}_0), \label{eq:shearing}
\end{align}
where $\mat{B} = \mathrm{diag}(EI_1, EI_2, GJ)$ is the bending/torsion
stiffness matrix, $\mat{S}$ the shear/stretch stiffness, $\bm{\kappa}$ the
curvature strain vector, and $\krest$ the \emph{rest curvature} that serves
as the control input.

The bending energy that drives actuation is:
\begin{equation}
  E_{\mathrm{bend}} = \int_0^L \frac{1}{2}\,
    (\bm{\kappa} - \krest)^\top \mat{B}\, (\bm{\kappa} - \krest) \, ds.
  \label{eq:bending-energy}
\end{equation}
Setting $\krest$ to a time-varying serpenoid profile makes the elastic
restoring forces produce undulatory locomotion.

% ===================================================================
\section{Two-Level Time-Stepping Hierarchy}
\label{sec:hierarchy}
% ===================================================================

\begin{table}[H]
\centering
\caption{Time-stepping levels and their default values.}
\label{tab:timesteps}
\renewcommand{\arraystretch}{1.3}
\begin{tabular}{@{}llll@{}}
\toprule
\textbf{Level} & \textbf{Symbol} & \textbf{Default} & \textbf{Description} \\
\midrule
Integration substep & $\dtsub$ & $0.001\;\mathrm{s}$ &
  Smallest timestep passed to the symplectic integrator \\
Intermediate block & $\Delta t$ & $0.05\;\mathrm{s}$ &
  $\Nsub = 50$ substeps; legacy config unit \\
RL action step & $\dtrl$ & $0.5\;\mathrm{s}$ &
  $\Nrl = 500$ substeps per agent decision \\
\bottomrule
\end{tabular}
\end{table}

The relationships are:
\begin{equation}
  \dtsub = \frac{\Delta t}{\Nsub}, \qquad
  \dtrl  = \Nrl \cdot \dtsub.
  \label{eq:dt-relations}
\end{equation}

At each RL step the agent outputs action parameters (amplitude, frequency,
wave number, phase, turn bias).  These parameters are held fixed for $\Nrl$
substeps, but the resulting curvature field evolves continuously because the
serpenoid curve depends on the running time $t$.

% ===================================================================
\section{Symplectic Integrators}
\label{sec:integrators}
% ===================================================================

PyElastica provides two symplectic integrators:

\subsection{Position Verlet (Default)}
\label{sec:position-verlet}

The St\"ormer--Verlet method is a second-order symplectic integrator.
For a system $\dot{\vect{q}} = \vect{v}$, $\dot{\vect{v}} = \vect{a}(\vect{q})$,
one full step from $t$ to $t + \dtsub$ is:
\begin{align}
  \vect{q}_{1/2} &\gets \vect{q}_n + \tfrac{1}{2}\dtsub\,\vect{v}_n, \label{eq:pv1} \\
  \vect{v}_{n+1} &\gets \vect{v}_n + \dtsub\,\vect{a}(\vect{q}_{1/2}), \label{eq:pv2} \\
  \vect{q}_{n+1} &\gets \vect{q}_{1/2} + \tfrac{1}{2}\dtsub\,\vect{v}_{n+1}. \label{eq:pv3}
\end{align}

\noindent This is time-reversible and preserves the symplectic structure,
meaning energy drift is bounded over long integrations.

\subsection{PEFRL (Position Extended Forest--Ruth--Like)}
\label{sec:pefrl}

A fourth-order symplectic integrator with optimized coefficients
\cite{omelyan2002optimized}.  It uses four force evaluations per step
(versus one for Verlet) but achieves much smaller energy error at the
same $\dtsub$.  The update sequence is:
\begin{align}
  \vect{q}_1 &\gets \vect{q}_n + \xi\,\dtsub\,\vect{v}_n, \notag \\
  \vect{v}_1 &\gets \vect{v}_n + (1-2\lambda)\tfrac{\dtsub}{2}\,\vect{a}(\vect{q}_1), \notag \\
  \vect{q}_2 &\gets \vect{q}_1 + \chi\,\dtsub\,\vect{v}_1, \notag \\
  \vect{v}_2 &\gets \vect{v}_1 + \lambda\,\dtsub\,\vect{a}(\vect{q}_2), \notag \\
  \vect{q}_3 &\gets \vect{q}_2 + (1-2(\chi+\xi))\,\dtsub\,\vect{v}_2, \notag \\
  \vect{v}_3 &\gets \vect{v}_2 + \lambda\,\dtsub\,\vect{a}(\vect{q}_3), \notag \\
  \vect{q}_4 &\gets \vect{q}_3 + \chi\,\dtsub\,\vect{v}_3, \notag \\
  \vect{v}_{n+1} &\gets \vect{v}_3 + (1-2\lambda)\tfrac{\dtsub}{2}\,\vect{a}(\vect{q}_4), \notag \\
  \vect{q}_{n+1} &\gets \vect{q}_4 + \xi\,\dtsub\,\vect{v}_{n+1},
  \label{eq:pefrl}
\end{align}
with $\xi \approx 0.1786$, $\lambda \approx -0.2123$, $\chi \approx -0.6626$.

% ===================================================================
\section{Stability Constraints}
\label{sec:stability}
% ===================================================================

Since both integrators are explicit, the substep $\dtsub$ must satisfy
a CFL-like condition dictated by the stiffest mode in the discretized rod.
The critical timescale is set by bending wave propagation:
\begin{equation}
  \dtsub < C \sqrt{\frac{\rho A \, \Delta s^4}{E I}},
  \label{eq:cfl}
\end{equation}
where $\Delta s = L / N_e$ is the element length, $N_e$ the number of
elements, and $C$ is an order-unity constant that depends on the
integrator (smaller for Verlet, larger for PEFRL).

\begin{table}[H]
\centering
\caption{Default material and geometry parameters.}
\label{tab:params}
\renewcommand{\arraystretch}{1.3}
\begin{tabular}{@{}lll@{}}
\toprule
\textbf{Parameter} & \textbf{Symbol} & \textbf{Value} \\
\midrule
Young's modulus     & $E$           & $10^5\;\mathrm{Pa}$ \\
Poisson ratio       & $\nu$         & $0.5$ \\
Shear modulus       & $G = E/2(1+\nu)$ & $3.33 \times 10^4\;\mathrm{Pa}$ \\
Density             & $\rho$        & $1200\;\mathrm{kg/m^3}$ \\
Rod length          & $L$           & $0.5\;\mathrm{m}$ \\
Rod radius          & $r$           & $0.02\;\mathrm{m}$ \\
Number of elements  & $N_e$         & $20$ \\
Element length      & $\Delta s$    & $0.025\;\mathrm{m}$ \\
Second moment of area & $I = \pi r^4 / 4$ & $1.257 \times 10^{-7}\;\mathrm{m^4}$ \\
\bottomrule
\end{tabular}
\end{table}

Substituting the default parameters into (\ref{eq:cfl}):
\begin{equation}
  \sqrt{\frac{\rho A \, \Delta s^4}{E I}}
  = \sqrt{\frac{1200 \times \pi (0.02)^2 \times (0.025)^4}
               {10^5 \times 1.257 \times 10^{-7}}}
  \approx 1.7 \times 10^{-3}\;\mathrm{s}.
\end{equation}
The default substep of $\dtsub = 10^{-3}\;\mathrm{s}$ sits safely below this
limit, providing a stability margin of roughly $1.7\times$.

\subsection{Numerical Damping}
\label{sec:damping}

PyElastica applies an \texttt{AnalyticalLinearDamper} with coefficient
$\gamma$.  This adds a dissipative term $-\gamma \vect{v}$ to each node's
velocity at every substep.

\begin{itemize}
  \item $\gamma = 0.002$: default for locomotion --- preserves undulatory
        dynamics while suppressing numerical ringing.
  \item $\gamma \geq 0.008$: \textbf{freezes the rod} completely; overdamped
        regime where elastic restoring forces cannot overcome dissipation.
\end{itemize}

\noindent The valid range is therefore $\gamma \in (0,\; 0.008)$.

% ===================================================================
\section{CPG--Integrator Synchronization}
\label{sec:cpg-sync}
% ===================================================================

The serpenoid curvature profile is:
\begin{equation}
  \kappa(s_j, t) = A \sin\!\bigl(k\, s_j + \omega\, t + \phi\bigr) + \kappa_{\mathrm{turn}},
  \label{eq:serpenoid}
\end{equation}
where $s_j$ are joint arc-length positions, $A$ is amplitude, $k$ the
spatial wave number, $\omega = 2\pi f$ the temporal frequency, $\phi$ a
phase offset, and $\kappa_{\mathrm{turn}}$ a constant steering bias.

A valid integration requires that $\krest$ in (\ref{eq:bending}) is updated
from (\ref{eq:serpenoid}) \emph{at every substep}, not just at RL decision
boundaries.  Between RL steps, the action parameters $(A, f, k, \phi,
\kappa_{\mathrm{turn}})$ are held constant but $t$ advances by $\dtsub$
each substep, so the curvature field evolves smoothly.

% ===================================================================
\section{Ground Friction at Each Substep}
\label{sec:friction}
% ===================================================================

Anisotropic Resistive Force Theory (RFT) friction is applied before each
integration substep:
\begin{equation}
  \vect{f}_{\mathrm{RFT}} = -c_t\,\vect{v}_t - c_n\,\vect{v}_n,
  \label{eq:rft}
\end{equation}
where $\vect{v}_t$ and $\vect{v}_n$ are the tangential and normal
components of node velocity, with default coefficients $c_t = 0.01$ and
$c_n = 0.05$ (a $5{:}1$ anisotropy ratio enabling serpentine propulsion).

The friction forces must be included in the external forcing term
$\vect{f}_{\mathrm{ext}}$ of (\ref{eq:linear-momentum}) at every substep.
Applying them only at RL boundaries would produce non-physical jerky
dynamics and destabilize the integrator.

% ===================================================================
\section{Contrast: DisMech Implicit Time Integration}
\label{sec:dismech}
% ===================================================================

The project supports a second physics backend, DisMech, which uses
\emph{implicit} (backward Euler) time integration.  Understanding the
differences clarifies why the validity conditions for Elastica exist.

\subsection{Backward Euler Formulation}
\label{sec:backward-euler}

DisMech discretizes the equations of motion (\ref{eq:linear-momentum})--(\ref{eq:angular-momentum})
using the backward Euler scheme.  Given the state vector
$\vect{x} = [\vect{q}^\top,\, \vect{v}^\top]^\top$ (positions and velocities
of all $N$ nodes), the update rule is:
\begin{align}
  \vect{v}_{n+1} &= \vect{v}_n + \Delta t\,\mat{M}^{-1}\vect{F}(\vect{q}_{n+1},\, \vect{v}_{n+1}),
    \label{eq:be-vel} \\
  \vect{q}_{n+1} &= \vect{q}_n + \Delta t\,\vect{v}_{n+1},
    \label{eq:be-pos}
\end{align}
where $\mat{M}$ is the mass matrix and $\vect{F}$ includes all internal
elastic forces, damping, and external forces.  Crucially, $\vect{F}$ is
evaluated at the \emph{unknown} future state $(\vect{q}_{n+1}, \vect{v}_{n+1})$,
making the system implicit.

Substituting (\ref{eq:be-pos}) into (\ref{eq:be-vel}) and rearranging gives
the nonlinear residual:
\begin{equation}
  \vect{g}(\vect{q}_{n+1})
    \coloneqq \mat{M}\bigl(\vect{q}_{n+1} - 2\vect{q}_n + \vect{q}_{n-1}\bigr)
    - \Delta t^2\,\vect{F}(\vect{q}_{n+1})
    = \vect{0},
  \label{eq:residual}
\end{equation}
which is solved via Newton's method at each timestep.

\subsection{Newton Solver}
\label{sec:newton}

Each timestep requires iterating:
\begin{equation}
  \vect{q}^{(k+1)} = \vect{q}^{(k)}
    - \mat{J}^{-1}(\vect{q}^{(k)})\,\vect{g}(\vect{q}^{(k)}),
  \label{eq:newton}
\end{equation}
where $\mat{J} = \partial\vect{g}/\partial\vect{q}$ is the Jacobian
(tangent stiffness matrix plus mass and damping contributions).
Convergence is declared when:
\begin{equation}
  \|\vect{g}(\vect{q}^{(k)})\| < \epsilon_f
  \quad\text{or}\quad
  \|\vect{q}^{(k+1)} - \vect{q}^{(k)}\| < \epsilon_d,
  \label{eq:convergence}
\end{equation}
with default tolerances $\epsilon_f = 10^{-4}$ (force) and
$\epsilon_d = 10^{-2}$ (displacement), and a maximum of $25$ iterations.

\subsection{Key Differences from PyElastica}
\label{sec:comparison}

\begin{table}[H]
\centering
\caption{Comparison of time integration approaches.}
\label{tab:comparison}
\renewcommand{\arraystretch}{1.3}
\begin{tabular}{@{}lll@{}}
\toprule
\textbf{Property} & \textbf{PyElastica (Explicit)} & \textbf{DisMech (Implicit)} \\
\midrule
Integration scheme
  & Symplectic (Verlet / PEFRL)
  & Backward Euler (Newton) \\
Timestep per call
  & $\dtsub = 0.001\;\mathrm{s}$
  & $\Delta t = 0.05\;\mathrm{s}$ \\
Steps per second
  & $1000$
  & $20$ \\
Stability
  & \emph{Conditional} ($\dtsub < \dtsub^{\max}$)
  & \emph{Unconditional} \\
Substep loop
  & $\Nrl = 500$ substeps per RL action
  & $1$ step per RL action \\
Cost per step
  & Low (explicit evaluation)
  & High (Jacobian assembly + linear solve) \\
Energy behavior
  & Bounded drift (symplectic)
  & Numerical dissipation (implicit) \\
Curvature control
  & \texttt{rest\_kappa} array
  & Natural strain of bend springs \\
Solver parameters
  & $\dtsub$, $\gamma$ (damping)
  & $\epsilon_f$, $\epsilon_d$, max\_iter \\
\bottomrule
\end{tabular}
\end{table}

\noindent The fundamental trade-off is:

\begin{itemize}
  \item \textbf{PyElastica} takes many cheap explicit steps.  Each substep
        is a simple algebraic update with no matrix inversion, but the
        timestep is constrained by the CFL condition~(\ref{eq:cfl}).
        The symplectic structure ensures long-term energy conservation.
  \item \textbf{DisMech} takes few expensive implicit steps.  Each step
        requires assembling and factorizing the Jacobian $\mat{J}$
        (typically sparse, solved via PARDISO in the C++ backend), but
        the method is unconditionally stable so $\Delta t$ can be $50\times$
        larger.  However, backward Euler introduces \emph{numerical
        dissipation} that damps high-frequency modes --- acceptable for
        locomotion but less faithful for energy-sensitive analyses.
\end{itemize}

\subsection{Why Validity Conditions Differ}
\label{sec:why-differ}

Because DisMech is unconditionally stable, many of the Elastica validity
conditions in Section~\ref{sec:summary} do not apply:

\begin{enumerate}[nosep]
  \item \textbf{No CFL constraint} --- DisMech can use $\Delta t = 0.05\;\mathrm{s}$
        directly, without substeps.
  \item \textbf{No per-substep curvature update} --- with only one step per
        RL action, curvature is set once at the step boundary (the CPG
        profile is evaluated at $t_n$, not continuously during the step).
  \item \textbf{No explicit damping coefficient} --- numerical dissipation
        from backward Euler itself provides damping; an additional
        damping viscosity is optional (default $0.0$).
  \item \textbf{Convergence monitoring replaces stability monitoring} ---
        the relevant diagnostic is the Newton residual norm
        $\|\vect{g}\|$, not the energy drift.
\end{enumerate}

\noindent This means DisMech is simpler to configure but offers less
temporal resolution of the CPG waveform within each RL step.  The
continuous curvature evolution that Elastica captures across $500$
substeps is collapsed to a single snapshot in DisMech.

% ===================================================================
\section{Algorithms}
\label{sec:algorithms}
% ===================================================================

% -------------------------------------------------------------------
\begin{algorithm}[t]
\caption{Single RL Step with PyElastica Integration}\label{alg:rl-step}
\begin{algorithmic}[1]
\Require Action $\vect{a} = (A,\, f,\, k,\, \phi,\, \kappa_{\mathrm{turn}})$
\Require Substep size $\dtsub$, substeps per action $\Nrl$
\Require Rod state $(\vect{q},\, \vect{v})$, simulation time $t$
\Ensure Updated state $(\vect{q}',\, \vect{v}')$, new time $t'$
\State $\omega \gets 2\pi f$ \Comment{Angular frequency}
\For{$i = 1$ \textbf{to} $\Nrl$}
    \For{$j = 1$ \textbf{to} $N_e - 1$} \Comment{Update rest curvature}
        \State $\kappa_{\mathrm{rest},j} \gets A \sin(k\, s_j + \omega\, t + \phi) + \kappa_{\mathrm{turn}}$
    \EndFor
    \State $\vect{f}_{\mathrm{ext}} \gets \Call{ComputeRFT}{\vect{v},\, c_t,\, c_n}$ \Comment{Friction}
    \State $(\vect{q},\, \vect{v},\, t) \gets \Call{SymplecticStep}{\vect{q},\, \vect{v},\, t,\, \dtsub}$
\EndFor
\State \textbf{return} $(\vect{q},\, \vect{v},\, t)$
\end{algorithmic}
\end{algorithm}
% -------------------------------------------------------------------

% -------------------------------------------------------------------
\begin{algorithm}[t]
\caption{Position Verlet Substep (\textsc{SymplecticStep})}\label{alg:verlet}
\begin{algorithmic}[1]
\Require Positions $\vect{q}$, velocities $\vect{v}$, time $t$, substep $\dtsub$
\Ensure Updated $(\vect{q}',\, \vect{v}',\, t')$
\State $\vect{q}_{1/2} \gets \vect{q} + \frac{1}{2}\dtsub\,\vect{v}$
    \Comment{Half-step position}
\State $\vect{a} \gets \vect{a}(\vect{q}_{1/2})$
    \Comment{Acceleration from internal + external forces}
\State $\vect{v}' \gets \vect{v} + \dtsub\,\vect{a}$
    \Comment{Full-step velocity}
\State $\vect{q}' \gets \vect{q}_{1/2} + \frac{1}{2}\dtsub\,\vect{v}'$
    \Comment{Half-step position}
\State $\vect{v}' \gets \vect{v}' - \gamma\,\vect{v}'$
    \Comment{Damping ($\gamma \ll 1$)}
\State $t' \gets t + \dtsub$
\State \textbf{return} $(\vect{q}',\, \vect{v}',\, t')$
\end{algorithmic}
\end{algorithm}
% -------------------------------------------------------------------

% -------------------------------------------------------------------
\begin{algorithm}[t]
\caption{Resistive Force Theory (\textsc{ComputeRFT})}\label{alg:rft}
\begin{algorithmic}[1]
\Require Node velocities $\vect{v} \in \Real^{3 \times N}$, drag coefficients $c_t, c_n$
\Ensure External friction forces $\vect{f} \in \Real^{3 \times N}$
\For{$i = 1$ \textbf{to} $N$}
    \State $\hat{\vect{t}}_i \gets$ unit tangent at node $i$
        \Comment{From adjacent positions}
    \State $v_{t,i} \gets (\vect{v}_i \cdot \hat{\vect{t}}_i)\,\hat{\vect{t}}_i$
        \Comment{Tangential component}
    \State $v_{n,i} \gets \vect{v}_i - v_{t,i}$
        \Comment{Normal component}
    \State $\vect{f}_i \gets -c_t\, v_{t,i} - c_n\, v_{n,i}$
\EndFor
\State \textbf{return} $\vect{f}$
\end{algorithmic}
\end{algorithm}
% -------------------------------------------------------------------

% -------------------------------------------------------------------
\begin{algorithm}[t]
\caption{Single RL Step with DisMech Integration (Implicit)}\label{alg:dismech-step}
\begin{algorithmic}[1]
\Require Action $\vect{a} = (A,\, f,\, k,\, \phi,\, \kappa_{\mathrm{turn}})$
\Require Timestep $\Delta t$, rod state $(\vect{q},\, \vect{v})$, time $t$
\Require Tolerances $\epsilon_f$, $\epsilon_d$, max iterations $K$
\Ensure Updated state $(\vect{q}',\, \vect{v}')$, residual norm $r$
\State $\omega \gets 2\pi f$
\For{$j = 1$ \textbf{to} $N_e - 1$} \Comment{Set curvature once per RL step}
    \State $\kappa_{\mathrm{nat},j} \gets A \sin(k\, s_j + \omega\, t + \phi) + \kappa_{\mathrm{turn}}$
\EndFor
\State $\vect{q}^{(0)} \gets \vect{q} + \Delta t\,\vect{v}$ \Comment{Initial guess (explicit Euler)}
\For{$k = 0$ \textbf{to} $K - 1$} \Comment{Newton iterations}
    \State $\vect{g} \gets \mat{M}(\vect{q}^{(k)} - 2\vect{q} + \vect{q}_{\mathrm{prev}}) - \Delta t^2\,\vect{F}(\vect{q}^{(k)})$
    \State $\mat{J} \gets \frac{\partial \vect{g}}{\partial \vect{q}}\big|_{\vect{q}^{(k)}}$
        \Comment{Jacobian (sparse)}
    \State $\delta\vect{q} \gets \mat{J}^{-1}\vect{g}$ \Comment{Sparse linear solve}
    \State $\vect{q}^{(k+1)} \gets \vect{q}^{(k)} - \delta\vect{q}$
    \If{$\|\vect{g}\| < \epsilon_f$ \textbf{or} $\|\delta\vect{q}\| < \epsilon_d$}
        \State \textbf{break} \Comment{Converged}
    \EndIf
\EndFor
\State $\vect{v}' \gets (\vect{q}^{(k+1)} - \vect{q}) / \Delta t$
\State $r \gets \|\vect{g}\|$ \Comment{Final residual for diagnostics}
\State \textbf{return} $(\vect{q}^{(k+1)},\, \vect{v}',\, r)$
\end{algorithmic}
\end{algorithm}
% -------------------------------------------------------------------

% ===================================================================
\section{Summary of Validity Conditions}
\label{sec:summary}
% ===================================================================

A time integration is \emph{valid} if and only if all conditions for the
chosen backend are satisfied.

\subsection{PyElastica (Explicit Symplectic)}

\begin{enumerate}
  \item \textbf{Substep stability:} $\dtsub < C\sqrt{\rho A \,\Delta s^4 / (EI)}$,
        satisfied by $\dtsub = 10^{-3}\;\mathrm{s}$ with the default parameters.

  \item \textbf{Symplectic integrator:} Either Position Verlet (second-order)
        or PEFRL (fourth-order) is used.  Non-symplectic integrators
        (e.g., forward Euler) would accumulate energy drift.

  \item \textbf{Per-substep curvature update:} The serpenoid
        $\kappa(s_j, t)$ is recomputed at every substep using the
        current simulation time $t$, not cached from the RL boundary.

  \item \textbf{Per-substep friction:} RFT (or Coulomb/Stribeck) forces
        are computed and applied before each integrator substep.

  \item \textbf{Bounded damping:} $\gamma \in (0,\; 0.008)$ to avoid
        freezing the rod while suppressing numerical ringing.

  \item \textbf{Consistent time accounting:} Both the integrator time
        and CPG phase advance by exactly $\dtsub$ per substep; no
        time desynchronization between actuation and physics.
\end{enumerate}

\subsection{DisMech (Implicit Backward Euler)}

\begin{enumerate}
  \item \textbf{Newton convergence:} The residual norm $\|\vect{g}\|$
        must fall below $\epsilon_f$ (or displacement below $\epsilon_d$)
        within the maximum iteration budget.  A diverging or stagnating
        Newton solve indicates an invalid step.

  \item \textbf{Timestep--material consistency:} While unconditionally
        stable, excessively large $\Delta t$ degrades accuracy through
        numerical dissipation.  The default $\Delta t = 0.05\;\mathrm{s}$
        is validated against the locomotion frequency band.

  \item \textbf{Curvature set before solve:} The natural strain of each
        bend spring must be set to the target curvature \emph{before}
        calling the Newton solver, so the elastic energy landscape
        reflects the desired actuation.

  \item \textbf{Residual monitoring:} The returned residual norm should
        be logged; values above $\epsilon_f$ indicate incomplete
        convergence and unreliable state.
\end{enumerate}

% ===================================================================
% References
% ===================================================================
\begin{thebibliography}{9}

\bibitem{gazzola2018forward}
M.~Gazzola, L.~H. Dudte, A.~G. McCormick, and L.~Mahadevan,
``Forward and inverse problems in the mechanics of soft filaments,''
\textit{Royal Society Open Science}, vol.~5, no.~6, 2018.

\bibitem{zhang2019pyelastica}
X.~Zhang, F.~K. Chan, T.~Parthasarathy, and M.~Gazzola,
``Modeling and simulation of complex dynamic musculoskeletal
architectures,''
\textit{Nature Communications}, vol.~10, 2019.

\bibitem{omelyan2002optimized}
I.~P. Omelyan, I.~M. Mryglod, and R.~Folk,
``Optimized Forest--Ruth- and Suzuki-like algorithms for integration
of motion in many-body systems,''
\textit{Computer Physics Communications}, vol.~146, pp.~188--199, 2002.

\bibitem{bergou2008discrete}
M.~Bergou, M.~Wardetzky, S.~Robinson, B.~Audoly, and E.~Grinspun,
``Discrete elastic rods,''
\textit{ACM Transactions on Graphics (SIGGRAPH)}, vol.~27, no.~3, 2008.

\bibitem{bergou2010discrete}
M.~Bergou, B.~Audoly, E.~Vouga, M.~Wardetzky, and E.~Grinspun,
``Discrete viscous threads,''
\textit{ACM Transactions on Graphics (SIGGRAPH)}, vol.~29, no.~4, 2010.

\end{thebibliography}

\end{document}
